\chapter{JavaScript}

Finora abbiamo imparato a costruire pagine web moderne e curate: con HTML ne
definiamo la struttura, con i CSS il loro aspetto sui vari media.

Le nostre pagine restano però \textbf{intrinsecamente statiche}: non c'è modo di
far cambiare il loro contenuto mentre le stiamo consultando, a meno di
modificare il file HTML e ricaricare la pagina. Manca il terzo pilastro.

\section{Che cos'è JavaScript}

\dfn{JavaScript}{
  \term{JavaScript} è un linguaggio di scripting di alto livello e a
  \textbf{tipizzazione debole}. Introdotto per la prima volta dal browser
  Netscape nel \textbf{1995}, è stato progettato per essere incluso nelle pagine
  web ed eseguito \textbf{dentro un browser}, allo scopo di rendere i documenti
  HTML più dinamici. I programmi JavaScript (gli \textit{script}) vengono
  interpretati come testo semplice dal \textit{motore JS}.}

Oggi JavaScript non si usa più soltanto nelle pagine web: lo si trova nel
software di back-end, nelle applicazioni desktop e mobili, nello scripting
generico. È il linguaggio di programmazione più diffuso in assoluto: nel 2023 il
66\% degli sviluppatori professionisti intervistati da Stack Overflow dichiarava
di usarlo.

JavaScript implementa la specifica \term{ECMAScript}.

\info{Perché due nomi}{
  \textit{ECMAScript} è lo \textbf{standard}, mantenuto dall'ente di
  normazione ECMA; \textit{JavaScript} è il nome comune delle sue
  implementazioni. Le sigle come ES5, ES6 (o ES2015) indicano le edizioni dello
  standard, ed è così che si datano le funzionalità del linguaggio: dire che una
  cosa è \guillemotleft ES6\guillemotright\ significa dire che è stata introdotta nel 2015.

  Il nome, va detto, è un artificio di marketing: JavaScript non ha nulla a che
  vedere con Java. La somiglianza sintattica è superficiale, la semantica è
  profondamente diversa.}

\section{Includere JavaScript in una pagina}

Il codice JavaScript si può includere in un documento HTML in due modi.

\textbf{JavaScript interno}: il codice è scritto dentro un elemento \tg{script},
nell'\tg{head} oppure nel \tg{body}.

\begin{lstlisting}[style=html]
<script>
  console.log("Hello World!");
</script>
\end{lstlisting}

\textbf{JavaScript esterno}: si usa \code{<script src="url/dello/script.js">}, e
il file esterno deve avere estensione \code{.js}.

\begin{lstlisting}[style=html]
<script src="script.js"></script>
\end{lstlisting}

\begin{esempio}{Hello World}
\begin{lstlisting}[style=html, name=hello.html]
<!DOCTYPE html>
<html>
  <head>
    <meta charset="utf-8">
    <title>JavaScript</title>
    <script>
      console.log("Hello World!");
    </script>
  </head>
  <body>
    <h1>Hello!</h1>
  </body>
</html>
\end{lstlisting}

\begin{console}
Hello World!
\end{console}

Il messaggio non compare nella pagina: compare nella \textbf{console} dei
DevTools. È il canale principale con cui uno script comunica a chi lo sviluppa.
\end{esempio}

\warning{Il tag script non è mai auto-chiudente}{
  Anche quando si carica un file esterno e il tag è vuoto, va scritto
  \code{<script src="..."></script>}. La forma abbreviata
  \code{<script src="..." />} non funziona: la specifica HTML non lo consente,
  e il browser considera tutto ciò che segue come contenuto dello script.}

\section{JavaScript moderno}

Per molto tempo JavaScript si è evoluto \textbf{senza rotture}: le nuove
funzionalità venivano aggiunte e quelle esistenti restavano intatte. Questo è
stato vero fino al 2009, con l'arrivo di \term{ECMAScript 5} (ES5), che ha
modificato alcune funzionalità esistenti.

Per mantenere la retro-compatibilità, queste modifiche incompatibili sono
\textbf{disattivate} per impostazione predefinita. Si abilitano con la direttiva
\code{"use strict"} nella prima riga di uno script.

\begin{lstlisting}[style=js, name=script.js]
"use strict";
/* qui il codice JavaScript */
\end{lstlisting}

In questa dispensa, salvo indicazione contraria, si fa sempre riferimento alla
versione moderna e \term{strict} di JavaScript.

\info{Perché conviene sempre}{
  La modalità strict non aggiunge funzionalità: \textbf{toglie} comportamenti
  permissivi che nascondono errori. Trasforma in errori espliciti cose che
  altrimenti fallirebbero in silenzio, come assegnare a una variabile mai
  dichiarata o scrivere su una proprietà in sola lettura (li vedremo entrambi
  più avanti in questo capitolo). Costa una riga e fa emergere subito
  bug che altrimenti si scoprirebbero molto più tardi.}

\section{Il linguaggio}

JavaScript supporta i paradigmi \textbf{imperativo}, \textbf{funzionale} e
\textbf{orientato agli oggetti}, e possiede costrutti molto interessanti per la
programmazione asincrona.

Un programma JavaScript è una \textbf{sequenza di istruzioni}
(\textit{statement}), composte da valori, operatori, espressioni, parole chiave
e commenti. La sintassi è piuttosto simile a quella di Java e del C, ma
\textbf{molto più permissiva}.

Le istruzioni sono delimitate da un a-capo, da un punto e virgola, o da
entrambi.

\begin{lstlisting}[style=js, name=istruzioni.js]
// con i punti e virgola
msg = "Hello";
num = 1;
console.log(msg);

// senza punti e virgola: funziona lo stesso
msg = "Hello"
num = 1
console.log(msg)

// piu' istruzioni sulla stessa riga
msg = "Hello"; num = 1; console.log(msg);
\end{lstlisting}

Delimitare le istruzioni con il punto e virgola resta comunque una buona
pratica da preferire.

\warning{L'inserimento automatico dei punti e virgola}{
  Quando il punto e virgola manca, l'interprete lo inserisce da sé
  (\textit{Automatic Semicolon Insertion}). Il meccanismo funziona quasi sempre,
  ma ha casi limite che generano bug memorabili: un \code{return} seguito da un
  a-capo, per esempio, restituisce \code{undefined} perché il punto e virgola
  viene inserito subito dopo il \code{return}. Metterli sempre costa poco.}

\section{Variabili}

\subsection{Il ciclo di vita di una variabile}

Dichiarare una variabile in JavaScript consiste di tre passi distinti.

\begin{figure}[H]
  \centering
  \begin{tikzpicture}[font=\Lato\small,
      st/.style={rounded corners=3pt, draw=primary!65!black, line width=0.9pt,
                 fill=primary!10!white, minimum width=3.4cm, minimum height=0.9cm,
                 align=center},
      a/.style={-{Stealth[length=6pt,width=5pt]}, line width=1pt,
                draw=neutral!45!black}]

    \node[st] (d) at (0,0)    {\textbf{Dichiarazione}};
    \node[st] (i) at (0,-2.1) {\textbf{Inizializzazione}};
    \node[st] (u) at (0,-4.2) {\textbf{Uso}};

    \draw[a] (d) -- (i);
    \draw[a] (i) -- (u);

    \node[font=\Lato\scriptsize, anchor=west, align=left, text=neutral!25!black]
      at (2.1,0) {il nome della variabile viene legato\\ allo scope corrente};
    \node[font=\Lato\scriptsize, anchor=west, align=left, text=neutral!25!black]
      at (2.1,-2.1) {la variabile viene inizializzata};
    \node[font=\Lato\scriptsize, anchor=west, align=left, text=neutral!25!black]
      at (2.1,-4.2) {la variabile può essere referenziata};

    \draw[decorate, decoration={brace, amplitude=5pt}, draw=red-gradient-5,
          line width=0.9pt] (-2.0,-1.75) -- (-2.0,-0.35);
    \node[font=\Lato\scriptsize, anchor=east, align=right, text=red-gradient-6!75!black]
      at (-2.25,-1.05) {\textbf{Temporal Dead Zone}\\ la variabile è\\ inaccessibile};
  \end{tikzpicture}
  \caption{Le tre fasi della vita di una variabile e la Temporal Dead Zone.}
  \label{fig:var-lifecycle}
\end{figure}

\subsection{let e const}

Nel JavaScript moderno le variabili si dichiarano e inizializzano con:

\begin{itemize}
  \item la parola chiave \code{let} per le variabili \guillemotleft normali\guillemotright;
  \item la parola chiave \code{const} per le costanti, che non possono essere
        riassegnate.
\end{itemize}

Per impostazione predefinita le variabili sono inizializzate a
\code{undefined}.

\begin{lstlisting}[style=js, name=variabili.js]
let x;
const y = 42;

console.log(x);   // output: undefined
console.log(y);   // output: 42
\end{lstlisting}

Con \code{const} si dichiarano \textbf{costanti locali con scope di blocco},
cioè variabili il cui valore non può essere riassegnato. Per il resto si
comportano come quelle dichiarate con \code{let}, compreso l'hoisting che
vedremo fra poco.

\begin{lstlisting}[style=js, name=const.js]
const n = 42;
console.log(n);   // output: 42
n = 43;           // TypeError: Assignment to constant variable.
\end{lstlisting}

\begin{warningbox}{{const non rende immutabile un oggetto}}
  \code{const} impedisce di \textbf{riassegnare la variabile}, non di
  modificare il valore a cui si riferisce. Se il valore è un oggetto, la
  variabile ne contiene solo il \textbf{riferimento} (lo vedremo nella sezione
  sugli oggetti), e le proprietà si possono cambiare liberamente.

\begin{lstlisting}[style=js]
const pet = { name: "Hannibal" };
pet.name = "Sif";   // consentito: cambia l'oggetto, non la variabile
pet.age = 7;        // consentito
pet = {};           // TypeError: e' la variabile a non poter cambiare
\end{lstlisting}
\end{warningbox}

\subsection{Gli scope}

Esistono tre livelli di scope:

\begin{itemize}
  \item \term{scope globale};
  \item \term{scope di funzione};
  \item \term{scope di modulo}.
\end{itemize}

In aggiunta, le variabili dichiarate con \code{let} e \code{const} possono avere
uno \term{scope di blocco}, cioè il più vicino blocco delimitato da una coppia
di parentesi graffe.

\begin{esempio}{Scope di blocco con let}
\begin{lstlisting}[style=js, name=scope.js]
let bool;              // dichiarazione
bool = true;

if (bool) {
  // console.log(msg); // msg non e' accessibile qui
  let msg = "Hello";   // dichiarazione + assegnamento
  console.log(msg);    // output: Hello
}

// console.log(msg);   // msg non e' definita qui
console.log(bool);     // output: true
\end{lstlisting}

Lo scope delle variabili dichiarate con \code{let} e \code{const} è il blocco
più vicino, e le variabili sono accessibili \textbf{solo dopo} che la riga in
cui sono dichiarate è stata eseguita.
\end{esempio}

\subsection{Hoisting}

\dfn{Hoisting}{
  L'\term{hoisting} (letteralmente \guillemotleft sollevamento\guillemotright) è il comportamento
  predefinito per cui le dichiarazioni di variabili e funzioni vengono
  implicitamente spostate all'\textbf{inizio del loro scope}.}

Con \code{let} e \code{const} viene sollevata soltanto la \textbf{dichiarazione},
non l'inizializzazione: quest'ultima avviene sulla riga che conteneva
originariamente la dichiarazione. Nell'intervallo fra le due la variabile esiste
ma è inaccessibile: è la \term{Temporal Dead Zone} della figura precedente.

\begin{lstlisting}[style=js, name=hoisting.js]
// qui la variabile x non e' definita
{
  // la dichiarazione di x viene sollevata qui: inizia la TDZ
  console.log(x);    // ReferenceError
  let x = 1;         // fine della TDZ, x viene inizializzata
  // da qui in poi x vale 1
}
// qui la variabile x non e' piu' definita
\end{lstlisting}

\begin{console}
Uncaught ReferenceError: can't access lexical declaration
'x' before initialization
\end{console}

\begin{esempio}{Shadowing}
\begin{lstlisting}[style=js, name=shadowing.js]
let n;                // dichiarazione
n = 1;

if (n > 0) {
  let n = 2;          // maschera la n dello scope esterno
  // let n;           // errore: non si puo' ridichiarare nello stesso blocco
  console.log(n);     // output: 2
}

console.log(n);       // output: 1
\end{lstlisting}
\end{esempio}

\subsection{Prima di ECMAScript 6: var e dichiarazioni implicite}

Prima dell'introduzione di ECMAScript 6 (2015) le variabili si dichiaravano con
la parola chiave \code{var}, o addirittura implicitamente. A meno di dover
davvero supportare browser molto vecchi, \textbf{non} vanno più dichiarate in
questi modi: hanno proprietà piuttosto confuse e favoriscono gli errori.

Le variabili dichiarate con \code{var} vengono sollevate allo scope di funzione
o globale più vicino (\textbf{niente scope di blocco}) e vengono \textbf{anche
inizializzate} a \code{undefined}.

\begin{lstlisting}[style=js, name=var.js]
if (true) {
  console.log(n);   // output: undefined (non un errore!)
  var n = 1;
}
console.log(n);     // output: 1 (n e' sopravvissuta al blocco)
\end{lstlisting}

Il codice sopra si comporta esattamente come questo:

\begin{lstlisting}[style=js, name=var-equivalente.js]
var n;
if (true) {
  console.log(n);   // output: undefined
  n = 1;
}
console.log(n);     // output: 1
\end{lstlisting}

Le variabili possono anche essere dichiarate \textbf{implicitamente}, per
esempio usandole in un assegnamento senza alcuna parola chiave. Il risultato è
la creazione di una \textbf{variabile globale}.

% Nota: qui serve la forma ambiente (e non la macro \error), perche' un
% listato verbatim non puo' comparire dentro l'argomento di una macro.
\begin{errorbox}{{Mai dichiarare implicitamente}}
  È una pessima pratica e una fonte inesauribile di errori: una variabile che
  doveva essere locale diventa globale, e da lì in poi qualunque altra parte del
  programma può sovrascriverla. In modalità \term{strict} la cosa è
  esplicitamente vietata e produce un \code{ReferenceError}.

\begin{lstlisting}[style=js]
"use strict";
msg = "meglio evitare";   // ReferenceError: msg is not defined
\end{lstlisting}
\end{errorbox}

\section{Tipi di dato primitivi}

L'operatore \code{typeof} restituisce il tipo di un valore.

\begin{lstlisting}[style=js, name=tipi.js]
let x;                        // dichiarata e inizializzata a undefined
console.log(typeof x);        // undefined

x = 42;
console.log(typeof x);        // number
x = 3.14;
console.log(typeof x);        // number
x = 10e12;
console.log(typeof x);        // number

x = "hello";
console.log(typeof x);        // string
x = 'hello';                  // apici singoli o doppi, indifferente
console.log(typeof x);        // string

x = false;
console.log(typeof x);        // boolean

x = 42 / 0;                   // Infinity
console.log(typeof x);        // number
x = 42 * "Hello";             // NaN
console.log(typeof x);        // number
\end{lstlisting}

\info{Un solo tipo numerico}{
  JavaScript non distingue interi e numeri in virgola mobile: esiste un solo
  tipo \code{number}, rappresentato in doppia precisione. Da qui due
  conseguenze da ricordare: \code{42/0} non è un errore ma vale
  \code{Infinity}, e il risultato di un'operazione numerica priva di senso è
  \code{NaN} (\textit{Not a Number}), che ha comunque tipo \code{number}.

  Curiosità utile all'esame: \code{NaN} è l'unico valore in JavaScript che non
  è uguale a se stesso. \code{NaN === NaN} vale \code{false}.}

\section{Operatori}

Gli operatori aritmetici sono quelli consueti.

\begin{center}
\renewcommand{\arraystretch}{1.35}
\begin{tabular}{@{}l@{\hspace{1.4cm}}l@{}}
  \textbf{Operatore} & \textbf{Descrizione} \\
  \code{=}  & assegnamento \\
  \code{+}  & addizione \\
  \code{-}  & sottrazione \\
  \code{*}  & moltiplicazione \\
  \code{**} & elevamento a potenza (da ECMAScript 2016) \\
  \code{/}  & divisione \\
  \code{\%} & modulo (resto della divisione) \\
  \code{++} & incremento \\
  \code{-{}-} & decremento \\
\end{tabular}
\end{center}

Gli operatori di confronto sono gli stessi di Java, con l'aggiunta di
\code{===} e \code{!==}.

\begin{center}
\renewcommand{\arraystretch}{1.35}
\begin{tabular}{@{}l@{\hspace{1.4cm}}l@{}}
  \textbf{Operatore} & \textbf{Descrizione} \\
  \code{==}  & uguale a \\
  \code{===} & stesso valore \textbf{e} stesso tipo \\
  \code{!=}  & diverso da \\
  \code{!==} & valore diverso \textbf{o} tipo diverso \\
  \code{>}   & maggiore di \\
  \code{<}   & minore di \\
  \code{>=}  & maggiore o uguale \\
  \code{<=}  & minore o uguale \\
  \code{?}   & operatore ternario \\
\end{tabular}
\end{center}

\subsection{Uguaglianza debole e uguaglianza stretta}

L'operatore di \term{uguaglianza debole} (\code{==}) verifica se i due operandi
sono uguali, \textbf{tentando di convertirli} quando sono di tipo diverso.

\begin{lstlisting}[style=js, name=uguaglianza-debole.js]
console.log(42 == 42);       // true
console.log("JS" == "JS");   // true
console.log("1" == 1);       // true
console.log(0 == false);     // true
console.log(true == 1);      // true
console.log(true == "1");    // true
console.log(true == 42);     // false
\end{lstlisting}

L'operatore di \term{uguaglianza stretta} (\code{===}) verifica l'uguaglianza
\textbf{senza} conversione di tipo: se gli operandi sono di tipo diverso, il
confronto restituisce immediatamente \code{false}.

\begin{lstlisting}[style=js, name=uguaglianza-stretta.js]
console.log(42 === 42);      // true
console.log("JS" === "JS");  // true
console.log("1" === 1);      // false
console.log(0 === false);    // false
console.log(true === 1);     // false
console.log(true === "1");   // false
console.log(true === 42);    // false
\end{lstlisting}

\warning{Usare sempre \code{===}}{
  L'ultima riga dei due esempi merita attenzione: \code{true == 1} è vero, ma
  \code{true == 42} è falso. La ragione è che la conversione trasforma
  \code{true} nel numero \code{1}, non in \guillemotleft un valore vero qualunque\guillemotright. Le
  regole di conversione dell'uguaglianza debole sono numerose e poco intuitive,
  e sono all'origine di una quantità sproporzionata di bug.

  La regola pratica è semplice: \textbf{usare sempre \code{===} e \code{!==}},
  e ricorrere a \code{==} solo quando si vuole davvero la conversione, sapendo
  esattamente che cosa si sta facendo.}

\section{Strutture di controllo}

\code{if}, \code{if-else}, \code{for}, \code{while}, \code{do}, \code{switch}
hanno la stessa sintassi e la stessa semantica di Java. Anche \code{continue} e
\code{break} funzionano esattamente come in Java.

\begin{lstlisting}[style=js, name=controllo.js]
if (condizione) {
  // codice
} else {
  // codice
}

for (let i = 0; i < 10; i++) {
  // codice
}

while (espressione) { /* codice */ }

do {
  // codice
} while (condizione);

switch (espressione) {
  case valore:
    // codice
    break;
  default:
    // codice
}
\end{lstlisting}

\section{Funzioni}

Le funzioni si dichiarano con la sintassi seguente.

\begin{lstlisting}[style=js, name=funzioni.js]
function greet(name) {
  console.log(`Hello ${name}`);   // backtick: template literal
}

greet("Web Technologies");        // stampa "Hello Web Technologies"
\end{lstlisting}

Lo scope di una dichiarazione di funzione è lo scope corrente, quello in cui la
dichiarazione avviene.

\info{I template literal}{
  Le stringhe delimitate da \term{backtick} (\code{`}) sono \textit{template
  literal}: al loro interno la sintassi \code{\$\{espressione\}} viene sostituita
  con il valore dell'espressione, e la stringa può occupare più righe. Sono
  molto più leggibili della concatenazione con \code{+}.

  Su tastiera italiana il backtick si inserisce con \code{ALT+096} dal
  tastierino numerico su Windows, con \code{Option+9} su macOS e con
  \code{AltGr+'} su Linux.}

\subsection{Argomenti predefiniti}

Quando in una chiamata non viene passato alcun argomento, i parametri sono
inizializzati a \code{undefined}. Le funzioni possono però dichiarare valori
predefiniti.

\begin{lstlisting}[style=js, name=default.js]
function greet(name, message = "Hello") {   // message ha un valore predefinito
  console.log(`${message} ${name}`);
}

greet("Web Technologies");           // Hello Web Technologies
greet("Web Technologies", "Ciao");   // Ciao Web Technologies
greet();                             // Hello undefined
\end{lstlisting}

\subsection{Hoisting delle funzioni}

L'hoisting si applica anche alle dichiarazioni di funzione, e in questo caso
solleva la funzione \textbf{per intero}: si può quindi invocare una funzione
\textbf{prima} della riga in cui è dichiarata.

\begin{lstlisting}[style=js, name=hoisting-funzioni.js]
greet("Web Technologies");   // stampa "Hello Web Technologies" (!)

function greet(name, message = "Hello") {
  console.log(`${message} ${name}`);
}
\end{lstlisting}

\subsection{Scope interno e visibilità}

Una funzione crea un proprio scope: le variabili e le funzioni dichiarate al suo
interno non sono visibili negli scope esterni. Al contrario, una funzione
\textbf{può accedere} alle variabili degli scope esterni (è una \term{closure}),
purché non siano mascherate nel proprio scope.

\begin{esempio}{Funzioni annidate e closure}
\begin{lstlisting}[style=js, name=closure.js]
let message = "Hello";

console.log(greet("Web Technologies"));  // output: Hello Web Technologies!

function greet(name) {                   // sollevata a inizio scope globale
  return generateMessage();

  function generateMessage() {           // sollevata a inizio scope di greet
    return `${message} ${name}!`;
  }
}

generateMessage();  // ReferenceError: generateMessage is not defined
\end{lstlisting}

La funzione \code{generateMessage} è locale allo scope di \code{greet} e non
può essere referenziata dall'esterno. Al suo interno, però, accede sia a
\code{name} (parametro di \code{greet}) sia a \code{message} (variabile
globale).
\end{esempio}

\subsection{Espressioni di funzione e arrow function}

Le funzioni possono anche essere create tramite \term{espressioni di funzione} e
assegnate a una variabile.

\begin{lstlisting}[style=js, name=espressioni.js]
// espressione di funzione classica
let greet = function (name) {
  console.log(`Hello ${name}!`);
}

// alternativa, con una arrow function
let greet = (name) => {
  console.log(`Hello ${name}!`);
}

greet("Web Technologies");   // output: Hello Web Technologies!
\end{lstlisting}

In questi casi valgono le regole di hoisting delle \textbf{variabili}, non
quelle delle funzioni. La differenza è tutt'altro che teorica.

\begin{esempio}{Hoisting a confronto}
\begin{lstlisting}[style=js, name=hoisting-confronto.js]
greet();              // ReferenceError: greet non e' definita
{
  greet();            // output: Hello
  function greet() {
    console.log("Hello");
  }
}
\end{lstlisting}

\begin{lstlisting}[style=js, name=hoisting-confronto2.js]
salute();             // ReferenceError: salute non e' definita
{
  salute();           // ReferenceError: variabile non inizializzata (TDZ)
  let salute = function () {
    console.log("Howdy");
  }
}
\end{lstlisting}

Nel primo caso la \textbf{dichiarazione di funzione} viene sollevata per intero
dentro il blocco, quindi la chiamata interna funziona. Nel secondo la variabile
\code{salute} viene sollevata ma non inizializzata: la chiamata cade nella
Temporal Dead Zone.

La prima riga fallisce in \textbf{entrambi} i casi perché, in modalità strict,
anche una dichiarazione di funzione scritta dentro un blocco ha \textbf{scope di
blocco}: viene sollevata all'inizio del blocco, non fuori. Senza
\code{"use strict"} i browser, per retro-compatibilità, rendono \code{greet}
visibile anche fuori dal blocco, e la prima chiamata darebbe un errore diverso:
un \code{TypeError} (\textit{greet is not a function}).
\end{esempio}

\subsection{Funzioni annidate restituite}

Una funzione è \term{annidata} quando viene creata dentro un'altra funzione. Le
funzioni annidate possono essere \textbf{restituite} e usate fuori dalla
funzione originale, e ovunque vengano usate mantengono l'accesso al
\textbf{contesto esterno} della funzione che le ha create.

\begin{esempio}{Una fabbrica di funzioni}
\begin{lstlisting}[style=js, name=fabbrica.js]
function getGreeter(message) {
  let sep = ",";
  return function (name) {
    console.log(`${message}${sep} ${name}!`);
  }
}

let helloGreeter = getGreeter("Hello");
helloGreeter("Web");        // Hello, Web!

let howdyGreeter = getGreeter("Howdy");
howdyGreeter("JS");         // Howdy, JS!
\end{lstlisting}

Le due funzioni restituite sono indipendenti e ciascuna \guillemotleft ricorda\guillemotright\ i
propri \code{message} e \code{sep}, anche se \code{getGreeter} è terminata da
tempo. È il meccanismo delle closure, ed è alla base di gran parte del
JavaScript idiomatico.
\end{esempio}

\section{Oggetti}

\dfn{Oggetto}{
  In JavaScript un \term{oggetto} è un contenitore di dati nella forma
  \code{chiave: valore}.}

\begin{lstlisting}[style=js, name=oggetti.js]
let a = new Object();   // sintassi "costruttore di oggetto"
let b = {};             // sintassi "oggetto letterale", usata piu' spesso
\end{lstlisting}

Le proprietà si possono aggiungere al momento della creazione:

\begin{lstlisting}[style=js, name=pet.js]
let pet = {
  name: "Hannibal",   // alla chiave "name" associa il valore "Hannibal"
  age: 7              // alla chiave "age" associa il valore 7
}
\end{lstlisting}

Una proprietà ha una \textbf{chiave} (detta anche nome o identificatore) prima
dei due punti, e un \textbf{valore} alla loro destra. Le dichiarazioni di
proprietà sono separate da virgole.

\subsection{Accedere alle proprietà}

Le proprietà si leggono con la \term{notazione a punto}.

\begin{lstlisting}[style=js, name=accesso.js]
let pet = { name: "Hannibal", age: 7 }

console.log(pet.name);       // output: Hannibal
console.log(pet.age);        // output: 7
console.log(pet.nickname);   // output: undefined
\end{lstlisting}

Accedere a una proprietà inesistente non è un errore: restituisce
\code{undefined}.

\subsection{Aggiungere ed eliminare proprietà}

\begin{lstlisting}[style=js, name=modifica.js]
let pet = { name: "Hannibal", age: 7 }

pet.nickname = "the Affable";   // nuova proprieta'
delete pet.age;                 // elimina la proprieta'

console.log(pet.name);          // "Hannibal"
console.log(pet.age);           // undefined
console.log(pet.nickname);      // "the Affable"

console.log(pet);               // stampa l'oggetto in console
console.table(pet);             // modo alternativo, in forma tabellare
\end{lstlisting}

\subsection{Notazione a parentesi quadre}

Le chiavi possono contenere spazi, e in tal caso vanno lette con la
\term{notazione a parentesi quadre}. Fra le parentesi si può usare
un'\textbf{espressione} qualunque.

\begin{lstlisting}[style=js, name=quadre.js]
let dog = {
  name: "Sif",
  "is a good boy": true
}

console.log(dog["name"]);
console.log(dog["is a good boy"]);

// come chiave si puo' usare un'espressione
dog["is the " + "best boy"] = true;
console.table(dog);
\end{lstlisting}

\subsection{Gli oggetti sono riferimenti}

Una variabile a cui è assegnato un oggetto memorizza un \term{riferimento}
all'oggetto, non l'oggetto stesso.

\begin{esempio}{Due nomi per lo stesso oggetto}
\begin{lstlisting}[style=js, name=riferimenti.js]
let firstPet = {
  name: "Richard",
  species: "Lizard"
}

let secondPet = firstPet;
secondPet.name = "Chuck";
secondPet.species = "Duck";

console.log(firstPet);   // Object { name: "Chuck", species: "Duck" }
console.log(secondPet);  // Object { name: "Chuck", species: "Duck" }

console.log(firstPet == secondPet);    // true, stesso valore
console.log(firstPet === secondPet);   // true, entrambi riferimenti uguali
\end{lstlisting}

Modificando \code{secondPet} è cambiato anche \code{firstPet}: non sono due
oggetti, è un solo oggetto con due nomi.
\end{esempio}

\subsection{Clonare un oggetto}

Per creare davvero una copia bisogna \textbf{iterare} sulle proprietà e
copiarle una a una.

\begin{lstlisting}[style=js, name=clone.js]
function clone(object) {
  let copy = {};                    // nuovo oggetto
  for (let key in object) {
    copy[key] = object[key];
  }
  return copy;
}

let pet = { name: "Richard", species: "Lizard" }
let copy = clone(pet);

copy.name = "Chuck";
console.log(pet.name);    // Richard
console.log(copy.name);   // Chuck
\end{lstlisting}

Il valore di una proprietà, però, può essere a sua volta un oggetto: la funzione
appena scritta produce una \term{copia superficiale} (\textit{shallow clone}),
in cui gli oggetti interni non vengono clonati ma solo referenziati.

\begin{esempio}{Il limite della copia superficiale}
\begin{lstlisting}[style=js, name=shallow.js]
let pet = {
  name: "Garfield",
  species: "Cat",
  owner: { name: "John", age: 17 }
}

let copy = clone(pet);
copy.owner.name = "Liz";

console.log(pet.owner.name);    // Liz  (!)
console.log(copy.owner.name);   // Liz
\end{lstlisting}

Per ottenere una \term{copia profonda} (\textit{deep clone}), cioè per clonare
anche gli oggetti interni, serve la ricorsione.

\begin{lstlisting}[style=js, name=deep-clone.js]
function clone(object) {
  let copy = {};
  for (let key in object) {
    if (typeof object[key] === "object") {
      copy[key] = clone(object[key]);   // ricorsione
    } else {
      copy[key] = object[key];
    }
  }
  return copy;
}
\end{lstlisting}

Implementare i metodi di clonazione da zero ha valore didattico, ma non è
necessario farlo ogni volta: esistono alternative già pronte, come
\code{Object.assign()} (superficiale) e \code{structuredClone()} (profonda).
\end{esempio}

\warning{Due trappole della copia profonda fatta a mano}{
  La versione ricorsiva decide se scendere di livello con il test
  \code{typeof object[key] === "object"}, che ha due punti deboli.
  \begin{itemize}
    \item \code{typeof null} vale \code{"object"} (una storica stranezza del
          linguaggio): una proprietà che vale \code{null} diventa, nella copia,
          un oggetto vuoto \code{\{\}}.
    \item Anche gli array hanno tipo \code{"object"}: \code{[1, 2]} diventa
          l'oggetto \code{\{0: 1, 1: 2\}}, che non è più un array.
  \end{itemize}
  \code{structuredClone()} gestisce correttamente entrambi i casi.}

\subsection{Metodi}

Le proprietà di un oggetto possono anche essere funzioni. Una funzione che è
proprietà di un oggetto si chiama \term{metodo}. I modi seguenti di definire un
metodo sono tutti equivalenti.

\begin{lstlisting}[style=js, name=metodi.js]
// (1) funzione anonima assegnata alla proprieta'
let p = {
  name: "John",
  age: 17,
  greet: function () { console.log("Hi!"); }
}

// (2) funzione dichiarata a parte e referenziata
let p = {
  name: "John",
  greet: greet
}
function greet() { console.log("Hi!"); }

// (3) proprieta' aggiunta dopo la creazione
let p = { name: "John" }
p.greet = function () { console.log("Hi!"); }

p.greet();   // Hi!
\end{lstlisting}

Esiste anche una \term{sintassi abbreviata} per dichiarare i metodi in un
oggetto letterale.

\begin{lstlisting}[style=js, name=shorthand.js]
// versione classica              // versione abbreviata
let p = {                         let p = {
  name: "John",                     name: "John",
  greet: function () {              greet() {
    console.log("Hi!");               console.log("Hi!");
  }                                 }
}                                 }
\end{lstlisting}

\subsection{La parola chiave this}

È comune che i metodi di un oggetto facciano riferimento ad altre proprietà
dello stesso oggetto. Vi accedono tramite la parola chiave \code{this}.

\begin{lstlisting}[style=js, name=this.js]
let john = {
  name: "John",
  greet() {
    console.log(`Hi, I'm ${this.name}!`);
  }
}

john.greet();   // Hi, I'm John!
\end{lstlisting}

Il punto cruciale è che il valore di \code{this} è valutato \textbf{a tempo di
esecuzione}, in base al contesto della chiamata.

\begin{esempio}{Lo stesso codice, due valori di this}
\begin{lstlisting}[style=js, name=this-runtime.js]
let john = { name: "John" };
let will = { name: "Will" };

let greet = function () {
  console.log(`Hi, I'm ${this.name}!`);
}

// la stessa funzione viene assegnata come metodo a due oggetti diversi
john.greet = greet;
will.greet = greet;

john.greet();   // Hi, I'm John!
will.greet();   // Hi, I'm Will!
\end{lstlisting}
\end{esempio}

\begin{warningbox}{{Le arrow function non hanno this}}
  Se \code{this} compare dentro una arrow function, viene preso dal
  \textbf{contesto esterno}, non dall'oggetto su cui il metodo è invocato.

\begin{lstlisting}[style=js, name=this-arrow.js]
let john = { nick: "John" };
let will = { nick: "Will" };

let greet = () => {
  console.log(`Hi, I'm ${this.nick}!`);
}

john.greet = greet;
will.greet = greet;

john.greet();   // Hi, I'm undefined!
will.greet();   // Hi, I'm undefined!
\end{lstlisting}

  Non è un difetto: è precisamente il motivo per cui le arrow function sono
  comode come callback, dove si vuole conservare il \code{this} circostante. La
  regola pratica: \textbf{mai} usare una arrow function per definire un metodo
  che deve accedere all'oggetto tramite \code{this}.
\end{warningbox}

\subsection{Costruttori}

Finora abbiamo creato oggetti con la sintassi letterale. E se dovessimo creare
molti oggetti simili fra loro? Si usano le \term{funzioni costruttore} e la
parola chiave \code{new}.

I costruttori sono normali funzioni, con due convenzioni:

\begin{enumerate}
  \item il loro nome inizia con una \textbf{lettera maiuscola};
  \item vanno invocati \textbf{solo} con la parola chiave \code{new}.
\end{enumerate}

\begin{lstlisting}[style=js, name=costruttore.js]
function Pet(name, species) {
  this.name = name;
  this.species = species;
  this.age = undefined;
}

let rick  = new Pet("Richard", "Lizard");
let chuck = new Pet("Chuck", "Duck");
\end{lstlisting}

Quando una funzione viene eseguita con \code{new}:

\begin{enumerate}
  \item viene creato un nuovo oggetto vuoto e assegnato a \code{this};
  \item viene eseguito il corpo della funzione, che tipicamente modifica
        \code{this};
  \item viene restituito il valore di \code{this}.
\end{enumerate}

\subsection{Optional chaining}

Accedendo a una proprietà inesistente si ottiene \code{undefined}. Ma accedere a
una proprietà \textbf{di} una proprietà inesistente produce un errore.

\begin{lstlisting}[style=js, name=chaining.js]
let pet = { name: "Garfield" }

console.log(pet.species);       // undefined
console.log(pet.owner.name);    // TypeError!
\end{lstlisting}

\info{{TypeError, non ReferenceError}}{
  Nelle slide l'errore di \code{pet.owner.name} è indicato come
  \code{ReferenceError}, ma quello che si ottiene davvero è un
  \code{TypeError} (\textit{Cannot read properties of undefined}). La
  differenza conta: un \code{ReferenceError} nasce quando si usa un
  \textbf{nome} che non esiste, come una variabile mai dichiarata; qui invece
  \code{pet} esiste e \code{pet.owner} vale semplicemente \code{undefined}, e
  l'errore sta nel leggere una proprietà da un valore che non è un oggetto.}

In molti casi pratici preferiremmo ottenere \code{undefined} (che significa
semplicemente \guillemotleft il nome del proprietario non è noto\guillemotright) piuttosto che un
errore. Si potrebbe verificare esplicitamente che la proprietà sia definita
prima di accedere a quelle interne, ma è ripetitivo e poco elegante:

\begin{lstlisting}[style=js]
let ownerName = pet.owner ? pet.owner.name : undefined;
\end{lstlisting}

L'operatore di \term{optional chaining} (\code{?.}) risolve la situazione:
interrompe immediatamente la valutazione (fa \textit{corto circuito}) se la
parte a sinistra è \code{undefined} (o \code{null}), restituendo
\code{undefined}.

\begin{lstlisting}[style=js]
let ownerName = pet.owner?.name;
\end{lstlisting}

\begin{esempio}{Varianti dell'optional chaining}
Si può usare anche per invocare una funzione che potrebbe non esistere, o per
accedere a proprietà con la notazione a parentesi quadre.

\begin{lstlisting}[style=js, name=chaining-varianti.js]
let john = {
  name: "John",
  greet() { return "Hi!"; },
  address: { "street name": "Web Dev Blvd" }
}

let mike = { name: "Mike", age: 17 }

console.log( john.greet?.() );                    // Hi!
console.log( mike.greet?.() );                    // undefined
console.log( john.address?.['street name'] );     // Web Dev Blvd
console.log( mike.address?.['street name'] );     // undefined
\end{lstlisting}
\end{esempio}

\subsection{Configurare le proprietà}

Oltre al valore, le proprietà di un oggetto hanno tre attributi speciali,
chiamati \term{flag}:

\begin{itemize}
  \item \term{writable}: se \code{true}, il valore può essere modificato;
        altrimenti la proprietà è in sola lettura;
  \item \term{enumerable}: se \code{true}, la proprietà compare nei cicli;
  \item \term{configurable}: se \code{true}, la proprietà può essere eliminata e
        i suoi flag possono essere modificati.
\end{itemize}

Quando si crea una proprietà nel modo \guillemotleft tradizionale\guillemotright, tutti i flag valgono
\code{true}. Le proprietà si possono però definire anche con
\code{Object.defineProperty()}, e in tal caso tutti i flag valgono
\code{false} per impostazione predefinita.

\begin{esempio}{Proprietà con flag espliciti}
\begin{lstlisting}[style=js, name=flags.js]
let pet = { species: "cat" }

Object.defineProperty(pet, "name", {
  value: "Garfield",
  writable: false,
  enumerable: false,
  configurable: false
})

for (let key in pet) {
  console.log(`Key: ${key}`);   // stampa solo "Key: species"
}

pet.name = "Odie";     // TypeError: "name" is read-only
delete pet.species;    // funziona
delete pet.name;       // TypeError: property "name" is non-configurable
\end{lstlisting}
\end{esempio}

\info{Gli errori li dà la modalità strict}{
  I due \code{TypeError} dell'esempio si ottengono in modalità strict. Senza
  \code{"use strict"} l'assegnamento e la \code{delete} falliscono \textbf{in
  silenzio}: \code{pet.name} resta \code{"Garfield"} e \code{delete} si limita
  a restituire \code{false}. È proprio uno dei fallimenti silenziosi che la
  modalità strict trasforma in errori espliciti.}

\subsection{Getter e setter}

Esistono due tipi di proprietà di un oggetto:

\begin{itemize}
  \item le \term{proprietà dati}, che memorizzano un valore: sono quelle viste
        finora;
  \item le \term{proprietà accessorie}, cioè funzioni eseguite quando la
        proprietà viene letta o scritta.
\end{itemize}

Le proprietà accessorie sono rappresentate da un \term{getter} (invocato in
lettura) e da un \term{setter} (invocato in scrittura), che si dichiarano con le
parole chiave \code{get} e \code{set}.

\begin{esempio}{Una proprietà calcolata}
\begin{lstlisting}[style=js, name=getset.js]
let p = {
  first: "John",
  last: "Smith",

  get fullName() {
    return `${this.first} ${this.last}`;
  },

  set fullName(name) {
    [this.first, this.last] = name.split(" ");   // assegnamento destrutturante
  }
}

// fullName si usa come una proprieta', non come una funzione: niente "()"
console.log(p.fullName);      // John Smith

p.fullName = "Jane White";
console.log(p.fullName);      // Jane White

p.fullName();                 // TypeError: p.fullName is not a function
\end{lstlisting}

Dall'esterno non c'è alcuna differenza fra una proprietà dati e una proprietà
accessoria: è precisamente questo il punto. Si può trasformare una proprietà
normale in una calcolata senza toccare il codice che la usa.
\end{esempio}

\section{Riferimenti}

\begin{itemize}
  \item \textbf{The Modern JavaScript Tutorial} \\
        \urlx{https://javascript.info/} \\
        Parte 1: \textit{An Introduction}, \textit{JavaScript Fundamentals},
        \textit{Code Quality} (3.1--3.4), \textit{Objects: the basics}
        (4.1--4.6), \textit{Data types} (5.1--5.3),
        \textit{Advanced working with functions} (6.3).
  \item \textbf{Eloquent JavaScript} (3\textsuperscript{a} edizione), Marijn
        Haverbeke. \\
        \urlx{https://eloquentjavascript.net/} \\
        Capitoli: Introduzione, da 1 a 6.
  \item \textbf{Learn JavaScript} --- MDN web docs. \\
        \urlx{https://developer.mozilla.org/en-US/docs/Learn/JavaScript} \\
        Utile come riferimento rapido, contiene molti esempi.
\end{itemize}
